\section{Lecture 3: Bases, Dimension, and Expansion}

% Increases line spacing within matrices and arrays globally for this section
\renewcommand{\arraystretch}{1.4} 
\vspace{1em}

\textbf{Intuition:} A basis is the perfect skeleton of a vector space. It is a set of vectors that spans the entire space (so you can reach absolutely everywhere) but is also linearly independent (so there is absolutely no redundancy in your coordinates). Because of these two properties, a basis provides a perfect, unique coordinate system for any space, ensuring every single vector has exactly one "address."

\vspace{2em}
% =========================================================================
\subsection{Bases and Standard Bases}
\vspace{1em}

\noindent\textbf{Definition (Basis).} \\[0.5em]
Let $V$ be a vector space. A subset $\mathcal{B} = \{\vec{v}_1, \dots, \vec{v}_k\} \subseteq V$ is called a \textbf{basis} for $V$ if it satisfies two conditions:
\begin{enumerate}[itemsep=1.5em, topsep=1em]
    \item $\mathcal{B}$ is linearly independent (no redundant vectors).
    \item $\operatorname{Span}(\mathcal{B}) = V$ (enough vectors to build the whole space).
\end{enumerate}
\vspace{0.5em}
\textit{Convention:} The basis for the trivial subspace $\{\vec{0}\}$ is the empty set $\emptyset$.

\vspace{1.5em}

\noindent\textbf{Standard Bases for Canonical Spaces} \\[0.5em]
Just like the $x, y, z$ axes in geometry, our abstract vector spaces have "standard" bases that are the simplest to work with:
\begin{itemize}[itemsep=1.5em, topsep=1em]
    \item \textbf{Euclidean Space $\mathbb{R}^3$:} The standard basis consists of the unit coordinate vectors:
    \[
    \mathcal{B} = \{ \vec{e}_1, \vec{e}_2, \vec{e}_3 \} = \left\{ \begin{bmatrix} 1 \\ 0 \\ 0 \end{bmatrix}, \begin{bmatrix} 0 \\ 1 \\ 0 \end{bmatrix}, \begin{bmatrix} 0 \\ 0 \\ 1 \end{bmatrix} \right\}.
    \]
    
    \item \textbf{Matrix Space $M_{2\times 2}(\mathbb{R})$:} The standard basis consists of single-entry unit matrices, one for each position:
    \[
    \mathcal{B} = \left\{ 
    \begin{bmatrix} 1 & 0 \\ 0 & 0 \end{bmatrix}, 
    \begin{bmatrix} 0 & 1 \\ 0 & 0 \end{bmatrix}, 
    \begin{bmatrix} 0 & 0 \\ 1 & 0 \end{bmatrix}, 
    \begin{bmatrix} 0 & 0 \\ 0 & 1 \end{bmatrix} 
    \right\} = \{E_{11}, E_{12}, E_{21}, E_{22}\}.
    \]
    
    \item \textbf{Polynomial Space $P_2(\mathbb{R})$:} The standard basis is the monomial basis, separating the constant, linear, and quadratic terms: $\mathcal{B} = \{1, x, x^2\}$.
\end{itemize}

\vspace{1.5em}

\begin{proof}[Verification that $\{1, x, x^2\}$ is a Basis for $P_2(\mathbb{R})$]
To rigorously prove this is a basis, we must check both criteria in the definition.
\begin{enumerate}[itemsep=1.5em, topsep=1em]
    \item \textbf{Spanning:} Let $p(x) \in P_2(\mathbb{R})$ be an arbitrary polynomial. By definition, it has the form $p(x) = a + bx + cx^2$ where $a, b, c \in \mathbb{R}$. We can write this exactly as a linear combination of our basis elements:
    \[
    p(x) = a(1) + b(x) + c(x^2).
    \]
    Since any $p(x)$ can be constructed this way, the set spans the space: $P_2(\mathbb{R}) = \operatorname{Span}\{1, x, x^2\}$.
    
    \item \textbf{Linear Independence:} Suppose a linear combination equals the zero polynomial, $0(x)$, which evaluates to $0$ for all $x \in \mathbb{R}$:
    \[ c_1(1) + c_2(x) + c_3(x^2) = 0(x) \]
    Because this equation is an identity (it must hold true for \textit{any} real number $x$), we can strategically plug in values for $x$ to solve for the coefficients.
    
    Evaluating at $x = 0$ kills the $x$ and $x^2$ terms, giving $c_1 = 0$. 
    
    Evaluating at $x = 1$ and $x = -1$ yields a small system:
    \[
    c_2 + c_3 = 0 \quad \text{and} \quad -c_2 + c_3 = 0 \implies 2c_3 = 0 \implies c_3 = 0.
    \]
    Substituting $c_3 = 0$ back in gives $c_2 = 0$.
    
    Thus, $c_1 = c_2 = c_3 = 0$. Because the trivial solution is the only solution, $\{1, x, x^2\}$ is linearly independent.
\end{enumerate}
\vspace{0.5em}
Because $\{1, x, x^2\}$ spans $P_2(\mathbb{R})$ and is linearly independent, it is a basis for $P_2(\mathbb{R})$.
\end{proof}

\vspace{2em}

\noindent \textbf{Theorem (Removing Redundant Vectors):} \\[0.5em]
Let $V$ be a non-empty vector space and suppose $\vec{v}_1, \dots, \vec{v}_k \in V$. If a vector $\vec{v}_i$ is a linear combination of the others, it is essentially "dead weight." Removing it does not alter the span of the set:
\[ 
\operatorname{Span}\{\vec{v}_1, \dots, \vec{v}_k\}
=
\operatorname{Span}\{\vec{v}_1, \dots, \vec{v}_{i-1}, \vec{v}_{i+1}, \dots, \vec{v}_k\}.
\]

\vspace{2em}
% =========================================================================
\subsection*{Worked Examples: Finding a Basis}
\vspace{1em}

\noindent\textbf{Example 1:} \\[0.5em]
Let $U = \operatorname{Span}\left\{
\begin{bmatrix} 1 \\ 1 \\ -1 \\ 3 \end{bmatrix},
\begin{bmatrix} -2 \\ -2 \\ 2 \\ -6 \end{bmatrix},
\begin{bmatrix} 2 \\ 1 \\ 1 \\ 2 \end{bmatrix},
\begin{bmatrix} 3 \\ 2 \\ 0 \\ 5 \end{bmatrix}
\right\}$. Find a basis for $U$.

\vspace{1em}

\begin{proof}[Solution]
To find a basis, we must filter out the redundant vectors. We first form a matrix whose columns are the given vectors. By row reducing this matrix, we track the dependencies between the columns.
\[
\begin{bmatrix}
 1 & -2 & 2 & 3 \\
 1 & -2 & 1 & 2 \\
-1 &  2 & 1 & 0 \\
 3 & -6 & 2 & 5
\end{bmatrix}
\]

We perform row reduction to find an REF (and its reduced form, RREF). First, we use the leading $1$ in the first column to eliminate all entries below it:
\[
\xrightarrow{\substack{R_2 \leftarrow R_2 - R_1 \\ R_3 \leftarrow R_3 + R_1 \\ R_4 \leftarrow R_4 - 3R_1}}
\begin{bmatrix}
1 & -2 &  2 &  3 \\
0 &  0 & -1 & -1 \\
0 &  0 &  3 &  3 \\
0 &  0 & -4 & -4
\end{bmatrix}
\]

Notice that the second column does not have a pivot (a leading non-zero entry). We move to the third column to establish our next pivot by multiplying row 2 by $-1$:
\[
\xrightarrow{R_2 \leftarrow -R_2}
\begin{bmatrix}
1 & -2 & 2 & 3 \\
0 &  0 & 1 & 1 \\
0 &  0 & 3 & 3 \\
0 &  0 & -4 & -4
\end{bmatrix}
\]

Now, we use this new pivot in the third column to eliminate the entries below it:
\[
\xrightarrow{\substack{R_3 \leftarrow R_3 - 3R_2 \\ R_4 \leftarrow R_4 + 4R_2}}
\begin{bmatrix}
1 & -2 & 2 & 3 \\
0 &  0 & 1 & 1 \\
0 &  0 & 0 & 0 \\
0 &  0 & 0 & 0
\end{bmatrix}
\]

Finally, we place the matrix into full Reduced Row Echelon Form (RREF) by clearing the entry above the pivot in column 3:
\[
\xrightarrow{R_1 \leftarrow R_1 - 2R_2}
\begin{bmatrix}
1 & -2 & 0 & 1 \\
0 &  0 & 1 & 1 \\
0 &  0 & 0 & 0 \\
0 &  0 & 0 & 0
\end{bmatrix}
\]

The leading entries (the pivots) are located in \textbf{column 1} and \textbf{column 3}. This tells us that columns 2 and 4 are free columns, meaning they are linear combinations of columns 1 and 3 (they are redundant). To form our basis, we discard the redundant columns and keep the \textbf{original vectors} corresponding to the pivot columns:
\[
\mathcal{B} = \left\{
    \begin{bmatrix} 1 \\ 1 \\ -1 \\ 3 \end{bmatrix},
    \begin{bmatrix} 2 \\ 1 \\ 1 \\ 2 \end{bmatrix}
\right\}
\]
\end{proof}

\vspace{2em}

\noindent\textbf{Example 2:} \\[0.5em]
Let $U = \operatorname{Span}\{2 - x + x^2,\, 1 + x - 2x^2,\, 3 - x^2\}$. Find a basis for $U$.

\vspace{1em}

\begin{proof}[Solution]
We must check these polynomials for linear independence. Suppose $c_1, c_2, c_3 \in \mathbb{R}$ such that their linear combination equals the zero polynomial:
\[
c_1(2-x+x^2) + c_2(1+x-2x^2) + c_3(3-x^2) = 0.
\]

If we expand and group the terms by their powers of $x$ (constants together, $x$ terms together, $x^2$ terms together), we force their coefficients to equal 0. This gives us the following \textbf{Linear System:}
\[
\begin{alignedat}{4}
2c_1 & {}+{} &  c_2 & {}+{} & 3c_3 & {}={} 0 \quad (\text{constants}) \\[0.5em]
 -c_1 & {}+{} &  c_2 &       &      & {}={} 0 \quad (x \text{ terms}) \\[0.5em]
  c_1 & {}-{} & 2c_2 & {}-{} &  c_3 & {}={} 0 \quad (x^2 \text{ terms})
\end{alignedat}
\]

We extract the coefficients into a matrix and perform row reduction to find the pivots:
\[
\quad\xrightarrow{\text{Coefficient Matrix}}\quad
\begin{bmatrix}
 2 &  1 &  3 \\
-1 &  1 &  0 \\
 1 & -2 & -1
\end{bmatrix}
\xrightarrow{\text{REF}}
\begin{bmatrix}
1 & 0 & 1 \\
0 & 1 & 1 \\
0 & 0 & 0
\end{bmatrix}
\]

The leading entries (pivots) appear in the first and second columns. This means the third column represents a redundant polynomial. Thus, we extract the first two original polynomials to form our basis:
\[
\mathcal{B} = \{2 - x + x^2,\, 1 + x - 2x^2\}
\]
\end{proof}

\vspace{2em}
% =========================================================================
\subsection{Dimension of a Vector Space}
\vspace{1em}

If $B = \{\vec{v}_1, \dots, \vec{v}_k\}$ and $C = \{\vec{w}_1, \dots, \vec{w}_l\}$ are both bases for a vector space $V$, then it is a mathematical guarantee that $k = l$. A vector space always requires the exact same number of basis vectors, regardless of which specific vectors you choose.

\vspace{1.5em}

\noindent\textbf{Definition (Dimension).} \\[0.5em]
If $B = \{\vec{v}_1, \dots, \vec{v}_k\}$ is a basis for a vector space $V$, then the \textbf{dimension} of $V$, denoted by $\dim(V)$, is the cardinality (size) of $B$. Written as: 
\[
\dim(V) = k \quad \text{In particular, if } V = \{\vec{0}\}, \text{ then } \dim(V) = 0.
\]

\vspace{1.5em}

\noindent\textbf{Examples of Dimension:} 
\begin{enumerate}[itemsep=1.5em, topsep=1em]
    \item $\dim(\mathbb{R}^n) = n$ (requires exactly $n$ coordinate vectors).
    \item $\dim(M_{m \times n}(\mathbb{R})) = m \times n$ (requires one vector for every matrix entry).
    \item $\dim(P_n(\mathbb{R})) = n + 1$ (the $+1$ is for the constant term).
\end{enumerate}

\vspace{1.5em}

\noindent\textbf{Theorem:} \\[0.5em]
Let $V$ be a finite dimensional space and suppose that $\dim(V) = n$. Then:
\begin{enumerate}[itemsep=1.5em, topsep=1em]
    \item A set of \textbf{more than $n$} vectors must be linearly dependent. (You have too many vectors, creating guaranteed redundancy).
    \item A set of \textbf{fewer than $n$} vectors cannot span $V$. (You don't have enough directions to reach every corner of the space).
    \item A set of \textbf{exactly $n$} vectors in $V$ spans $V$ if and only if it is linearly independent. (If the count is perfectly $n$, passing one test automatically guarantees passing the other).
\end{enumerate}

\vspace{1.5em}

\noindent\textbf{Examples:}
\begin{itemize}[itemsep=1.5em, topsep=1em]
    \item Let $B = \{1, x, x^2, 2x^2, -3x + 1\} \subseteq P_2(\mathbb{R})$. The set $B$ is linearly dependent because it contains 5 elements, which is more than $\dim(P_2(\mathbb{R})) = 3$. Because $5 > 3$, there is unavoidable redundancy. Therefore, $B$ cannot be a basis.
    
    \item Let $B = \left\{ \begin{bmatrix} 1 \\ 0 \\ 0 \end{bmatrix}, \begin{bmatrix} 0 \\ 1 \\ 0 \end{bmatrix} \right\} \subseteq \mathbb{R}^3$. Because there are only 2 vectors and $\dim(\mathbb{R}^3) = 3$, this set physically cannot span $\mathbb{R}^3$. It only describes a 2D plane.
\end{itemize}

\vspace{2em}
% =========================================================================
\subsection{Basis Expansion Theorem}
\vspace{1em}

\noindent\textbf{Theorem (Basis Expansion):} \\[0.5em]
Let $V$ be an $n$-dimensional vector space. If $C = \{\vec{v}_1, \dots, \vec{v}_k\} \subseteq V$ is a linearly independent set with $k < n$, then we can find additional vectors $\vec{v}_{k+1}, \dots, \vec{v}_n \in V$ such that:
\[
\{\vec{v}_1, \dots, \vec{v}_k, \vec{v}_{k+1}, \dots, \vec{v}_n\}
\]
is a full basis for $V$. (In other words, you can always "pad" an independent set until it becomes a basis).

\vspace{1.5em}

\noindent\textbf{Steps to expand a basis:}
\begin{enumerate}[itemsep=1.5em, topsep=1em]
    \item Consider the independent set.
    \item Form a matrix with a generic vector and find an REF to determine the dependency conditions.
    \item Choose a new vector that intentionally violates those conditions.
\end{enumerate}

\vspace{1.5em}

\noindent\textbf{Example 1 (Expanding in $\mathbb{R}^3$):} \\[0.5em]
Extend the set $C = \left\{ \begin{bmatrix} 1 \\ 2 \\ 3 \end{bmatrix}, \begin{bmatrix} 0 \\ 1 \\ 0 \end{bmatrix} \right\}$ to a basis for $\mathbb{R}^3$.

\vspace{1em}

\begin{proof}[Solution]
We note that $C$ is linearly independent (they are not scalar multiples). Since $\dim(\mathbb{R}^3) = 3$, we need to add exactly one more vector to $C$, ensuring it is NOT a linear combination of the vectors currently in $C$.

Let $\begin{bmatrix} a \\ b \\ c \end{bmatrix} \in \mathbb{R}^3$ be a generic vector. We check the condition for this vector to fall inside the span by checking if there exist $c_1, c_2 \in \mathbb{R}$ such that:
\[
\begin{bmatrix} a \\ b \\ c \end{bmatrix} = c_1 \begin{bmatrix} 1 \\ 2 \\ 3 \end{bmatrix} + c_2 \begin{bmatrix} 0 \\ 1 \\ 0 \end{bmatrix}
\]

Setting up the augmented matrix of this system and computing the REF, we get:
\[
\left(\begin{array}{cc|c} 
1 & 0 & a \\ 
0 & 1 & b - 2a \\ 
0 & 0 & c - 3a 
\end{array}\right)
\]

Looking at the bottom row, this system is inconsistent (meaning the vector is outside the span) when $c - 3a \neq 0 \implies c \neq 3a$. Choosing any numbers for $a, b, c$ that deliberately break this rule will give us a valid extension vector. 

Consider $a=2, b=0, c=1$. Because $1 \neq 3(2)$, the rule is broken, so the vector $\begin{bmatrix} 2 \\ 0 \\ 1 \end{bmatrix}$ is completely independent from the original two.

So, expanding the set gives:
\[
\mathcal{B} = \left\{ \begin{bmatrix} 1 \\ 2 \\ 3 \end{bmatrix}, \begin{bmatrix} 0 \\ 1 \\ 0 \end{bmatrix}, \begin{bmatrix} 2 \\ 0 \\ 1 \end{bmatrix} \right\}
\]
which is a full basis for $\mathbb{R}^3$.
\end{proof}

\vspace{2em}

\noindent\textbf{Example 2 (Expanding in Polynomial Space):} \\[0.5em]
Extend the set $C = \{x, 1+x\}$ to a basis for $P_3(\mathbb{R})$.

\vspace{1em}

\begin{proof}[Solution]
The space $P_3(\mathbb{R})$ is 4-dimensional, so we need two more polynomials that are not in the span of $C$. Notice that $C$ contains only degree 1 polynomials. Any combination of $C$ will only yield degree 1 or constant terms; it physically cannot create any polynomials of degree 2 or higher. 

Because of this, we can safely and simply add the standard higher-degree polynomials $x^2$ and $x^3$. Since they cannot be built from $C$, they maintain linear independence:
\[ \mathcal{B} = \{x, 1+x, x^2, x^3\} \]
\end{proof}

\end{document}